% GENERATED FILE. Edit canonical lessons, then run npm run generate:books.
% canonical-source-hash: fnv1a64:cc09e8faf328d533
% canonical-lessons: ES-C567-combate, ES-C567-tregua, ES-C567-batalla, ES-C567-soldado, ES-C567-ejercito, ES-R567-repaso-560-567

\chapter{La guerra}
\label{ch:es-la-guerra}
\hlchaptermodality{567}

\begin{chapteropening}
\textbf{By the end of this chapter:} \emph{I can talk about a war: the fighting, a truce, battles, soldiers and the army.}

Complete the last lesson of chapter 567: I can talk about a war: the fighting, a truce, battles, soldiers and the army.
\end{chapteropening}

\section[el combate]{el combate — the fight, the combat}

\label{lesson:ES-C567-combate}



\begin{quote}
Before the new one: say the Spanish for the record, then the Spanish for the novelty.
\end{quote}

\subsection*{You'll want to know: el combate}
\textbf{el combate} — \textquotedblleft{}the fight, the combat\textquotedblright{}.

\begin{quote}
Murieron cien hombres en el combate.
\end{quote}

\begin{quote}
\emph{A hundred men died in the fighting.}
\end{quote}

\emph{Fuera de combate} — out of action, knocked out.

\begin{cousinweb}
\emph{combate} is built on \emph{combatir} 'to fight', which comes from spoken Latin \emph{*combattere} (EtymDB), built on \emph{com-} and \emph{battuere} 'to beat'.
\end{cousinweb}

\begin{grammarlens}[title={What you've built}]
One more word for telling, explaining and describing at B1.
\end{grammarlens}

\subsection*{Your turn}
\begin{itemize}
\raggedright
  \item \emph{Say it:} \emph{el combate}
  \item \emph{Say it:} \emph{el combate}, once more
  \item \emph{Say it:} a sentence of your own with \emph{el combate}
  \item \emph{Recall:} say the Spanish for the record, then the Spanish for the novelty, then say \emph{el combate} again
\end{itemize}

\begin{tcolorbox}[breakable,colback=teal!4,colframe=teal!35!black,title={Before you move on}]
What does el combate mean? (\textquotedblleft{}The fight, the combat\textquotedblright{}.) Say it once more.
\end{tcolorbox}
\section[la tregua]{la tregua — the truce}

\label{lesson:ES-C567-tregua}



\begin{quote}
Before the new one: say the Spanish for the novelty, then the Spanish for the fight.
\end{quote}

\subsection*{You'll want to know: la tregua}
\textbf{la tregua} — \textquotedblleft{}the truce\textquotedblright{}.

\begin{quote}
Los dos países firmaron una tregua.
\end{quote}

\begin{quote}
\emph{The two countries signed a truce.}
\end{quote}

Also a break in anything hard: \emph{La lluvia no dio tregua.}

\begin{cousinweb}
\emph{tregua} comes from Gothic \emph{triggwa} 'a treaty, a covenant' (EtymDB), a Germanic word of the same family as English \emph{true} and \emph{truce}.
\end{cousinweb}

\begin{grammarlens}[title={What you've built}]
One more word for telling, explaining and describing at B1.
\end{grammarlens}

\subsection*{Your turn}
\begin{itemize}
\raggedright
  \item \emph{Say it:} \emph{la tregua}
  \item \emph{Say it:} \emph{la tregua}, once more
  \item \emph{Say it:} a sentence of your own with \emph{la tregua}
  \item \emph{Recall:} say the Spanish for the novelty, then the Spanish for the fight, then say \emph{la tregua} again
\end{itemize}

\begin{tcolorbox}[breakable,colback=teal!4,colframe=teal!35!black,title={Before you move on}]
What does la tregua mean? (\textquotedblleft{}The truce\textquotedblright{}.) Say it once more.
\end{tcolorbox}
\section[la batalla]{la batalla — the battle}

\label{lesson:ES-C567-batalla}



\begin{quote}
Before the new one: say the Spanish for the fight, then the Spanish for the truce.
\end{quote}

\subsection*{You'll want to know: la batalla}
\textbf{la batalla} — \textquotedblleft{}the battle\textquotedblright{}.

\begin{quote}
La batalla duró tres días.
\end{quote}

\begin{quote}
\emph{The battle lasted three days.}
\end{quote}

\begin{cousinweb}
\emph{batalla} comes from Late Latin \emph{battuālia} 'fighting and fencing drills' (EtymDB), from \emph{battuere} 'to beat', probably through Occitan.
\end{cousinweb}

\begin{grammarlens}[title={What you've built}]
One more word for telling, explaining and describing at B1.
\end{grammarlens}

\subsection*{Your turn}
\begin{itemize}
\raggedright
  \item \emph{Say it:} \emph{la batalla}
  \item \emph{Say it:} \emph{la batalla}, once more
  \item \emph{Say it:} a sentence of your own with \emph{la batalla}
  \item \emph{Recall:} say the Spanish for the fight, then the Spanish for the truce, then say \emph{la batalla} again
\end{itemize}

\begin{tcolorbox}[breakable,colback=teal!4,colframe=teal!35!black,title={Before you move on}]
What does la batalla mean? (\textquotedblleft{}The battle\textquotedblright{}.) Say it once more.
\end{tcolorbox}
\section[el soldado]{el soldado — the soldier}

\label{lesson:ES-C567-soldado}



\begin{quote}
Before the new one: say the Spanish for the truce, then the Spanish for the battle.
\end{quote}

\subsection*{You'll want to know: el soldado}
\textbf{el soldado} — \textquotedblleft{}the soldier\textquotedblright{}.

\begin{quote}
Los soldados volvieron a casa.
\end{quote}

\begin{quote}
\emph{The soldiers came home.}
\end{quote}

A woman is \emph{la soldado}.

\begin{cousinweb}
\emph{soldado} is usually traced to Italian \emph{soldato} 'paid', from \emph{soldo} 'pay', from Latin \emph{solidus}, a gold coin; EtymDB lists the two words as cousins. A soldier was a man paid to fight.
\end{cousinweb}

\begin{grammarlens}[title={What you've built}]
One more word for telling, explaining and describing at B1.
\end{grammarlens}

\subsection*{Your turn}
\begin{itemize}
\raggedright
  \item \emph{Say it:} \emph{el soldado}
  \item \emph{Say it:} \emph{el soldado}, once more
  \item \emph{Say it:} a sentence of your own with \emph{el soldado}
  \item \emph{Recall:} say the Spanish for the truce, then the Spanish for the battle, then say \emph{el soldado} again
\end{itemize}

\begin{tcolorbox}[breakable,colback=teal!4,colframe=teal!35!black,title={Before you move on}]
What does el soldado mean? (\textquotedblleft{}The soldier\textquotedblright{}.) Say it once more.
\end{tcolorbox}
\section[el ejército]{el ejército — the army}

\label{lesson:ES-C567-ejercito}



\begin{quote}
Before the new one: say the Spanish for the battle, then the Spanish for the soldier.
\end{quote}

\subsection*{You'll want to know: el ejército}
\textbf{el ejército} — \textquotedblleft{}the army\textquotedblright{}.

\begin{quote}
Mi tío estuvo en el ejército.
\end{quote}

\begin{quote}
\emph{My uncle was in the army.}
\end{quote}

\begin{cousinweb}
\emph{ejército} is a learned borrowing from Latin \emph{exercitus} 'a trained army' (EtymDB), from \emph{exercēre} 'to train, to exercise'.
\end{cousinweb}

\begin{grammarlens}[title={What you've built}]
One more word for telling, explaining and describing at B1.
\end{grammarlens}

\subsection*{Your turn}
\begin{itemize}
\raggedright
  \item \emph{Say it:} \emph{el ejército}
  \item \emph{Say it:} \emph{el ejército}, once more
  \item \emph{Say it:} a sentence of your own with \emph{el ejército}
  \item \emph{Recall:} say the Spanish for the battle, then the Spanish for the soldier, then say \emph{el ejército} again
\end{itemize}

\begin{tcolorbox}[breakable,colback=teal!4,colframe=teal!35!black,title={Before you move on}]
What does el ejército mean? (\textquotedblleft{}The army\textquotedblright{}.) Say it once more.
\end{tcolorbox}
\section[(dialogue)]{Repaso de los capítulos 560–567}

\label{lesson:ES-R567-repaso-560-567}



\begin{quote}
Say the Spanish for to gather and to get together.
\end{quote}

\subsection*{Your turn}
Say these aloud:

\begin{itemize}
\raggedright
  \item to gather, to get together, to call, the get-together, the reunion
  \item to leave, to set sail, the departure, the return, the absence
  \item to do again, to end up, to start, to burst out, to move away
  \item the fable, the account, to narrate, to relate, the narrator
  \item the character, the main character, the hero, the legend, the myth
  \item the plot, the storyline, the chapter, the episode, the ending
  \item the anecdote, the adventure, the feat, the record, the novelty
  \item the fight, the truce, the battle, the soldier, the army
\end{itemize}

\begin{tcolorbox}[breakable,colback=teal!4,colframe=teal!35!black,title={Before you move on}]
To gather? (\textbf{congregarse}.) To get together? (\textbf{juntarse}.) To call? (\textbf{convocar}.) The get-together? (\textbf{la tertulia}.) The reunion? (\textbf{el reencuentro}.) To leave? (\textbf{marcharse}.) To set sail? (\textbf{zarpar}.) The departure? (\textbf{la partida}.) The return? (\textbf{el regreso}.) The absence? (\textbf{la ausencia}.) To do again? (\textbf{volver a}.) To end up? (\textbf{acabar por}.) To start? (\textbf{ponerse a}.) To burst out? (\textbf{echarse a}.) To move away? (\textbf{alejarse}.) The fable? (\textbf{la fábula}.) The account? (\textbf{el relato}.) To narrate? (\textbf{narrar}.) To relate? (\textbf{relatar}.) The narrator? (\textbf{el narrador}.) The character? (\textbf{el personaje}.) The main character? (\textbf{el protagonista}.) The hero? (\textbf{el héroe}.) The legend? (\textbf{la leyenda}.) The myth? (\textbf{el mito}.) The plot? (\textbf{la trama}.) The storyline? (\textbf{el argumento}.) The chapter? (\textbf{el capítulo}.) The episode? (\textbf{el episodio}.) The ending? (\textbf{el desenlace}.) The anecdote? (\textbf{la anécdota}.) The adventure? (\textbf{la aventura}.) The feat? (\textbf{la hazaña}.) The record? (\textbf{el récord}.) The novelty? (\textbf{la novedad}.) The fight? (\textbf{el combate}.) The truce? (\textbf{la tregua}.) The battle? (\textbf{la batalla}.) The soldier? (\textbf{el soldado}.) The army? (\textbf{el ejército}.)
\end{tcolorbox}
